% Section 4.2, from lectures/L04/appendix/b_structures.md.

\section{Structures, Arrays, and Where They Live}
\label{c4:sec:structures}

\subsection{A structure is an agreement}
\label{c4:sec:agreement}
There is no \code{struct} here. There is an address, some bytes after it, and an agreement between
your code and itself about which byte means what.

\bookfigure{struct_layout}{An LED structure as seven byte cells at \code{0x0200} to \code{0x0206},
holding \code{0x23 0x00}, \code{0x24 0x00}, \code{0x25 0x00} and \code{0x05}, braced into the four
fields \code{pin_reg}, \code{dir_reg}, \code{port_reg} and \code{pin}, each labelled with the
\code{ldd} that reaches it.}{c4:fig:struct}

That is Chapter~2's LED structure after \code{led_init} has run for Arduino pin 13. Four things in
it are worth saying out loud, because in C the compiler would say them for you:

\textbf{A 16-bit field is two bytes, low byte first.} \code{0x23} followed by \code{0x00} is the
address \code{0x0023}, which is \code{PINB}. Read the two bytes in the other order and you get
\code{0x2300}, which is not memory at all, and the driver will store through it without complaint
(\secref{c4:sec:where}).

\textbf{The offsets are the interface.} They are declared in \file{drivers/include/led.inc} for
your driver, and copied into \file{device/include/avr/drivers.hpp} for the test suite. A structure
whose fields are in different places is a different structure, however similar the code looks.

\textbf{Every field is one instruction away.} All four offsets are below 63, so \code{ldd} and
\code{std} reach any of them directly (\secref{c4:sec:displacement}). That is the whole design: a
driver copies the structure address into \code{Z} once and then never does pointer arithmetic
again.

\textbf{Nothing checks the type.} A button structure begins with an LED structure's three
pointers, at the same offsets, so \code{led_on} will accept one and find three plausible port
registers in it. The fourth field is where they part: offset 6 is the LED's \code{pin} and the
button's \code{pcmsk_reg}, so \code{led_on} reads a pin number of \code{0x6B}, shifts by that, gets
a zero mask, writes nothing and reports success. There is no type system here; there is only what
you passed, and the failure is silence rather than an error.

\subsection{An array is a stride}
\label{c4:sec:stride}
An array of structures is a run of them, one after another, with nothing in between. To get from
one to the next you add the structure's \textbf{size}, and on this device that means:
\begin{itemize}
  \item the LED structure's stride is \textbf{7},
  \item the button structure's stride is \textbf{10}.
\end{itemize}
Neither is one, and neither is a power of two, and both of those facts have consequences.

\textbf{Adding 7 is not incrementing.} Walking an array is \code{adiw} by the structure size, or
the \code{subi}/\code{sbci} pair when the size is over 63. It is not \code{Z+}, which advances by
one byte and would put the second structure's first field on top of the first structure's second
field.

\textbf{Indexing needs a multiply.} The address of entry $n$ is the base plus $n$ times the stride,
and $7 \times n$ is not a shift. The AVR does have a hardware \code{mul}, which costs two cycles
and leaves its result in \code{r1:r0}; using it means remembering to clear \code{r1} afterwards,
because compiled C assumes \code{r1} holds zero (\secref{c2:sec:contract}). The alternative is to
walk from the start, which costs a loop and no surprises.

\textbf{A wrong stride is not an error.} It is a second structure that overlaps the first, and the
symptom is one LED that works and a second that behaves like a corrupted copy of it. There is
nothing to catch it: every address involved is a legal address, and every store succeeds.

\subsection{Passing an array}
\label{c4:sec:passing}
A subroutine that works on an array needs three things, and this book passes them in the argument
registers the calling contract gives it: the \textbf{base address} in \code{r25:r24}, the
\textbf{count} in \code{r22}, and, where there is one, an \textbf{index} in \code{r20}.

The count is not optional and it is not a courtesy. Nothing in memory marks where an array ends;
the byte after the last structure is an ordinary byte that will read as a plausible structure. A
subroutine that walks until something looks wrong walks forever.

\textbf{So every loop over an array tests its count at the top, not the bottom.} A count of zero
is a real input, it means ``do nothing'', and a loop that tests at the bottom does one entry
anyway. That is the single most common array bug, it is invisible for every other count, and the
test suite has a case for exactly it.

\subsection{Where a structure may live}
\label{c4:sec:where}

\bookfigure{sram_map}{A map of SRAM with high addresses at the top: not memory above
\code{0x0900}, the stack growing down from \code{RAMEND} at \code{0x08FF}, free space, then the
data segment reserved by \code{.byte} up from \code{0x0100}.}{c4:fig:sram}

An ATmega328P has 2048 bytes of SRAM, from \code{0x0100} to \code{0x08FF}. Above and below that are
addresses that a store instruction will accept and that are not memory you may use:

\begin{booktable}{@{}p{8.6em}p{8em}L@{}}
  \thead{Address} & \thead{What it is} & \thead{What a store does} \\
  \midrule
  \code{0x0000} to \code{0x001F} & the register file & changes a register, quietly \\
  \code{0x0020} to \code{0x005F} & the I/O registers & reconfigures the hardware \\
  \code{0x0060} to \code{0x00FF} & extended I/O & the same \\
  \code{0x0100} to \code{0x08FF} & SRAM & what you wanted \\
  \code{0x0900} upwards & \textbf{nothing} & reaches no memory; nothing reliable reads back \\
\end{booktable}

\textbf{\code{RAMEND + 1} is the one to know about.} It is \code{0x0900}, it looks like the obvious
place to start putting things, and a good deal of published AVR code allocates driver structures
there, including the material this book is based on. On a device with external RAM it \emph{is}
where that RAM begins. On an ATmega328P on an Arduino Uno there is no external RAM, so it is nothing
at all: the stores are accepted and reach no memory, the datasheet does not say what a load from
there returns, and the driver ends up working from values it never wrote. On a part that reads
them back as zeros, that looks as though every structure were full of zeros. The simulator is stricter than the chip here: simavr stops the
program at the first access above \code{RAMEND} and reports an invalid address, so under the test
suite this mistake is loud rather than quiet.

The right answer is to let the assembler place them. AVRASM2 has a data segment for exactly this:
\code{.dseg} switches to it, \code{label: .byte n} reserves $n$ bytes and gives you the address the
assembler chose, and \code{.cseg} switches back. Two reservations cannot overlap, and neither
occupies a byte of flash, because nothing in the data segment is emitted: it is an address
allocator and nothing more.

\begin{asm}
.dseg
my_led:     .byte LED_SIZE
my_button:  .byte BTN_SIZE
.cseg
\end{asm}

\textbf{Nothing zeroes it for you.} \code{.bss} in a C program is cleared by startup code that runs
before \code{main}, and there is no startup code here: \code{avra} has no linker and no runtime, so
the bytes hold whatever the SRAM powered up with. Every field a driver reads before it has written
is a field you initialised or a field that is garbage. That is what \code{led_init} is for.

Hard-coding \code{0x0200} works too and is what the test suite does, because a test wants to know
exactly where it put something; in a program, letting the assembler do it is one fewer number to
keep consistent.

\textbf{Whichever you choose, the arithmetic in \secref{c4:sec:stack} is still yours to do.}
Nothing on this device notices when your variables and your stack meet.
